% VI. Експериментални резултати и анализ.
% ВНИМАНИЕ: в този раздел не се пише нито едно число, преди то да е измерено.
\section{Експериментални резултати и анализ}
\label{sec:results}

Разделът представя получените резултати в реда на експериментите от
Раздел~\ref{sec:method}. Всяко число тук е измерено на машината, описана в
Табл.~\ref{tab:machine}, с командите, посочени в Раздел~\ref{sec:method}, и се намира в
суровите файлове в \code{results/}. Числа, които не са измерени, не са попълнени.

\subsection{Капацитет спрямо брой реплики}

Табл.~\ref{tab:throughput} съдържа и трите повторения на всяка точка, а не само
обобщение \tabref{tab:throughput}. Разсейването между повторенията е част от резултата и
скриването му би направило таблицата по-подредена и по-невярна.

\begin{table}[H]
\caption{Пропускателна способност и закъснение по персентили спрямо броя реплики. Затворен контур, 2 едновременни заявки на реплика, 5 s загряване, 20 s измерване}
\label{tab:throughput}
\centering
\begin{tabular}{@{}rrrrrrrr@{}}
\toprule
\textbf{Реплики} & \textbf{Повт.} & \textbf{Едновр.} & \textbf{Пропуск.} &
\textbf{p50} & \textbf{p90} & \textbf{p95} & \textbf{p99} \\
 & & & [зая./s] & [ms] & [ms] & [ms] & [ms] \\
\midrule
1 & 1 & 2  & 16{,}75 & 99{,}01  & 169{,}25 & 255{,}31 & 411{,}99 \\
1 & 2 & 2  & 19{,}80 & 99{,}25  & 106{,}20 & 109{,}44 & 120{,}55 \\
1 & 3 & 2  & 18{,}90 & 100{,}15 & 115{,}88 & 128{,}00 & 203{,}35 \\
\midrule
2 & 1 & 4  & 33{,}45 & 117{,}52 & 175{,}72 & 207{,}39 & 285{,}80 \\
2 & 2 & 4  & 21{,}45 & 179{,}52 & 282{,}37 & 321{,}73 & 433{,}96 \\
2 & 3 & 4  & 34{,}10 & 129{,}18 & 183{,}66 & 202{,}74 & 244{,}84 \\
\midrule
4 & 1 & 8  & 35{,}15 & 161{,}06 & 493{,}66 & 611{,}11 & 711{,}22 \\
4 & 2 & 8  & 59{,}90 & 104{,}38 & 232{,}10 & 304{,}19 & 462{,}00 \\
4 & 3 & 8  & 57{,}95 & 118{,}35 & 222{,}94 & 299{,}56 & 399{,}33 \\
\midrule
6 & 1 & 12 & 65{,}10 & 115{,}53 & 390{,}38 & 523{,}00 & 859{,}78 \\
6 & 2 & 12 & 77{,}60 & 120{,}22 & 244{,}28 & 353{,}46 & 888{,}83 \\
6 & 3 & 12 & 38{,}90 & 234{,}45 & 617{,}65 & 743{,}56 & 863{,}21 \\
\bottomrule
\end{tabular}
\end{table}

Неуспешни заявки няма: и в дванадесетте повторения делът им е нула.

Обобщението по медиана е дадено в Табл.~\ref{tab:throughput-median} \tabref{tab:throughput-median}.

\begin{table}[H]
\caption{Медиана от трите повторения и ускорение спрямо една реплика}
\label{tab:throughput-median}
\centering
\begin{tabular}{@{}rrrrrr@{}}
\toprule
\textbf{Реплики} & \textbf{Пропуск.} & \textbf{Ускорение} &
\textbf{p50} & \textbf{p90} & \textbf{p99} \\
 & [зая./s] & & [ms] & [ms] & [ms] \\
\midrule
1 & 18{,}90 & 1{,}00 & 99{,}25  & 115{,}88 & 203{,}35 \\
2 & 33{,}45 & 1{,}77 & 129{,}18 & 183{,}66 & 285{,}80 \\
4 & 57{,}95 & 3{,}07 & 118{,}35 & 232{,}10 & 462{,}00 \\
6 & 65{,}10 & 3{,}44 & 120{,}22 & 390{,}38 & 863{,}21 \\
\bottomrule
\end{tabular}
\end{table}

\begin{figure}[H]
\centering
\begin{tikzpicture}[x=1.55cm, y=0.055cm]
  % оси
  \draw[->, thick] (0,0) -- (7.2,0) node[right,font=\small] {реплики};
  \draw[->, thick] (0,0) -- (0,100) node[above,font=\small] {заявки/s};
  \foreach \x in {1,2,3,4,5,6} \draw (\x,0) -- (\x,-2) node[below,font=\scriptsize] {\x};
  \foreach \y in {0,20,40,60,80} \draw (0,\y) -- (-0.08,\y) node[left,font=\scriptsize] {\y};

  % идеално линейно мащабиране от измерената точка при една реплика
  \draw[dashed, gray] (1,18.90) -- (6,113.4);
  \node[gray, font=\scriptsize, anchor=west] at (4.3,96) {линейно};

  % измерена медиана
  \draw[very thick] (1,18.90) -- (2,33.45) -- (4,57.95) -- (6,65.10);
  \foreach \p in {(1,18.90),(2,33.45),(4,57.95),(6,65.10)} \fill \p circle (1.4pt);
  \node[font=\scriptsize, anchor=north west] at (6,65.10) {измерено};

  % разсейване между повторенията
  \foreach \x/\lo/\hi in {1/16.75/19.80, 2/21.45/34.10, 4/35.15/59.90, 6/38.90/77.60} {
    \draw[thick] (\x,\lo) -- (\x,\hi);
    \draw[thick] ({\x-0.05},\lo) -- ({\x+0.05},\lo);
    \draw[thick] ({\x-0.05},\hi) -- ({\x+0.05},\hi);
  }
\end{tikzpicture}
\caption{Пропускателна способност спрямо броя реплики. Точките са медианата от три повторения, отсечките показват най-малката и най-голямата стойност}
\label{fig:throughput}
\end{figure}

\textbf{Коментар.} Мащабирането е подлинейно и се изчерпва \figref{fig:throughput}. От 1
до 4 реплики ускорението е 3{,}07 при четирикратно увеличение на ресурса, тоест
ефективността е около 77\%. От 4 до 6 реплики пропускателната способност нараства само с
12\%, а закъснението на 99-и персентил се удвоява, от 462 ms на 863 ms. Обяснението не е
в услугата: при 6 реплики върху машина с 12 логически процесора едновременно работят 6
изчислителни нишки, 64 нишки на посредника и 12 нишки на генератора, и системата е
презаписана. Измерва се планировчикът на операционната система, не капацитетът на
услугата.

Разсейването нараства заедно с броя реплики. При една реплика разликата между най-ниската
и най-високата стойност е 3{,}05 заявки/s, при шест реплики е 38{,}70 заявки/s, тоест
почти колкото цялата медиана при две реплики. Този резултат е и основната причина
локалните числа да не се представят като характеристика на услугата: те са характеристика
на конкретната машина под конкретно натоварване.

Закъснението на 50-и персентил остава почти постоянно, между 99 и 130 ms, докато
закъснението на високите персентили расте многократно. Това е типичното поведение на
система с опашки и е точно причината в Раздел~\ref{sec:method} да е записано, че се
отчитат персентили, а не средна стойност: средното закъснение при 6 реплики би скрило
865 ms опашка зад 120 ms медиана.

\subsection{Време за реакция при увеличаване на набора}

\begin{table}[H]
\caption{Време от подаване на командата за увеличаване до насочване на трафик към новата реплика, 5 повторения}
\label{tab:reaction}
\centering
\begin{tabular}{@{}rrrr@{}}
\toprule
\textbf{Повторение} & \textbf{От} & \textbf{До} & \textbf{Време [s]} \\
\midrule
1 & 1 & 2 & 4{,}662 \\
2 & 1 & 2 & 5{,}234 \\
3 & 1 & 2 & 1{,}357 \\
4 & 1 & 2 & 2{,}001 \\
5 & 1 & 2 & 1{,}840 \\
\midrule
\multicolumn{3}{@{}l}{медиана} & 2{,}001 \\
\multicolumn{3}{@{}l}{най-малка, най-голяма} & 1{,}357 \dots 5{,}234 \\
\bottomrule
\end{tabular}
\end{table}

\textbf{Коментар.} Измереното време съдържа три съставки: създаване и стартиране на
контейнера, регистрация на новия адрес във вградения DNS, и откриването му от посредника,
чийто период на опресняване е 1 секунда. Разделянето им не е направено, затова се
съобщава общата стойност. Разликата от близо четири пъти между най-бързото и най-бавното
повторение при иначе еднакви условия показва, че времето за създаване на контейнер не е
постоянна величина, и че всяка политика за мащабиране, чийто период на управление е от
порядъка на секунда, работи с шум от същия порядък. Именно затова периодът на управление
в Експеримент 3 е 3 секунди, а не по-малък.

\subsection{Поведение на контролера под натоварване}

Стекът е вдигнат с една реплика, контролерът е стартиран с цел 0{,}70, граници от 1 до 6
реплики, период 3 s и периоди на изчакване 6 s при увеличаване и 15 s при намаляване.
Натоварването е в отворен контур, 40 заявки в секунда, 32 едновременни, 45 s измерване.

Отчетът на генератора за този период: 1868 успешни заявки, 0 неуспешни, пропускателна
способност 41{,}51 заявки/s, закъснение p50 = 51{,}84 ms, p90 = 1508{,}20 ms,
p95 = 1534{,}61 ms, p99 = 1639{,}47 ms, най-голямо 1658{,}75 ms. Разликата между медианата
и високите персентили е близо тридесеткратна и не е дефект: тя е опашката, натрупана в
първите секунди, когато една реплика приема натоварване, което не може да поеме. Това е
желаният вид резултат в отворен контур. Закъснението показва натрупаното чакане, вместо
клиентът тихо да намали скоростта си, а именно натрупаното чакане е онова, което
контролерът съществува да съкрати.

Табл.~\ref{tab:timeline} възпроизвежда времевата редица от решения \tabref{tab:timeline}.

\begin{table}[H]
\caption{Решения на контролера и метриката, която е предизвикала всяко от тях}
\label{tab:timeline}
\centering
\begin{tabular}{@{}rrrrL{5.4cm}@{}}
\toprule
\textbf{$t$ [s]} & \textbf{Реплики} & \textbf{Натоварване} & \textbf{Желани} & \textbf{Решение} \\
\midrule
3{,}0  & 1 & 0{,}000 & 1 & задържане, на целта \\
6{,}9  & 1 & 0{,}000 & 1 & задържане, на целта \\
10{,}3 & 1 & 1{,}101 & 2 & \textbf{увеличаване}, над целта \\
15{,}6 & 2 & 0{,}848 & 2 & задържане, изчакване преди увеличаване \\
20{,}2 & 2 & 0{,}822 & 3 & \textbf{увеличаване}, над целта \\
25{,}6 & 3 & 0{,}996 & 3 & задържане, изчакване преди увеличаване \\
30{,}2 & 3 & 0{,}531 & 3 & задържане, на целта \\
33{,}3 & 3 & 0{,}667 & 3 & задържане, на целта \\
36{,}3 & 3 & 0{,}654 & 3 & задържане, на целта \\
39{,}4 & 3 & 0{,}653 & 3 & задържане, на целта \\
42{,}5 & 3 & 0{,}661 & 3 & задържане, на целта \\
45{,}6 & 3 & 0{,}660 & 3 & задържане, на целта \\
48{,}7 & 3 & 0{,}654 & 3 & задържане, на целта \\
51{,}7 & 3 & 0{,}660 & 3 & задържане, на целта \\
54{,}8 & 3 & 0{,}655 & 3 & задържане, на целта \\
57{,}9 & 3 & 0{,}646 & 3 & задържане, на целта \\
60{,}9 & 3 & 0{,}033 & 1 & \textbf{намаляване}, под целта \\
65{,}9 & 1 & \multicolumn{2}{l}{нулиран брояч} & пропуснат интервал \\
68{,}9 & 1 & 0{,}000 & 1 & задържане, на целта \\
\multicolumn{5}{@{}l}{\dots\ следват още девет интервала на задържане до $t = 99{,}8$} \\
\bottomrule
\end{tabular}
\end{table}

\textbf{Коментар.} Контурът се затваря. Натоварването расте, контролерът увеличава набора
на две стъпки, от 1 на 2 и от 2 на 3 реплики, след което натоварването се установява
между 0{,}646 и 0{,}667, тоест непосредствено под целта 0{,}70, и контролерът спира да
действа за десет последователни интервала. След премахването на натоварването той връща
набора на една реплика.

Установената стойност позволява независима проверка на цялата верига. При 3 реплики,
натоварване 0{,}655 и измерени 40{,}0 заявки в секунда, времето за обслужване на една
заявка следва да е $0{,}655 \cdot 3 / 40{,}0 = 49{,}1$ ms. Това съвпада с времето, което
изчислителното ядро отчита само за себе си при същите параметри. Тоест метриката, която
управлява контролера, измерва действително извършената работа, а не сумата на чакането по
пътя до нея.

Три подробности заслужават внимание, защото са точно нещата, които политиката е
предвидена да прави.

Първо, при $t = 15{,}6$ натоварването 0{,}848 е над целта, но периодът на изчакване след
предишната промяна не е изтекъл и решението е задържане. Същото се случва при
$t = 25{,}6$ със стойност 0{,}996. Без периодите на изчакване контролерът щеше да добави
реплики на всеки интервал, докато първите още се стартират, и да надхвърли многократно
нужния брой, преди изобщо да види резултата от първото си действие.

Второ, при $t = 60{,}9$ натоварването пада на 0{,}033 и решението е намаляване направо до
долната граница. Скокът от 3 на 1 реплика не е грешка: желаният брой при това натоварване
е под единица, а долната граница е 1.

Трето, редът при $t = 65{,}9$ показва пропуснат интервал заради нулиран брояч. Общата
метрика на набора е сумата от броячите на репликите, така че при премахване на реплика тя
спада. Първата версия на контролера четеше отрицателната разлика като отрицателно
натоварване и намаляваше набора още повече. Обработката на нулирането е описана в
Раздел~\ref{sec:implementation}.

Зоната на нечувствителност не се задейства в това конкретно изпълнение, тъй като
установеното натоварване попада точно в интервала, при който желаният брой съвпада с
текущия. Поведението ѝ се проверява от модулен тест, а не от това измерване, и точно
затова политиката е отделена в чиста функция.

\subsection{Размер на контейнерния образ}

Измереният размер на получения образ е 8{,}09 MB, отчетен с \code{docker image ls}. Той
съдържа пет статично свързани изпълними файла и нищо друго: няма мениджър на пакети, няма
обвивка и няма споделена библиотека.

\subsection{Цена за единица извършена работа}

\TODO{Таблица със структура: доставчик, компонент на разхода, начислена сума за периода,
брой изпълнени единици работа, цена за единица. Не е попълнена, защото не съществува
начислена от доставчик сума: нищо не е разгръщано. Инструментът \code{stratus-cost}
изчислява таблицата при подадени цени и измерена пропускателна способност, но в него няма
вградена цена и такава няма да бъде добавена.}

\TODO{Коментар: коя среда излиза по-евтина за същата работа, кой компонент отговаря за
разликата, и променя ли се изводът при по-голям обем. Изисква Експеримент 4.}

\subsection{Сравнение между двамата доставчици}

\TODO{Таблица със структура: доставчик, брой заявки в секунда, закъснение на 50-и, 90-и,
95-и и 99-и персентил, дял на неуспешните заявки. Не е попълнена: описанието на
инфраструктурата не е прилагано, вж. Раздел~\ref{sec:method-scope}.}

\TODO{Таблица със структура: доставчик, време от преминаването на прага до създаване на
нова реплика, време до готовност, време до поемане на трафик, общо време за реакция. Не е
попълнена по същата причина.}

\subsection{Обобщен анализ}

Върху въпроса, поставен в увода, локалните резултати позволяват да се каже следното.

Построената система работи и контурът се затваря: метрика, изнесена от приложението,
управлява броя реплики, и решенията са проследими до числото, което ги е предизвикало
\tabref{tab:timeline}. Политиката е чиста функция и се проверява без работещ клъстер,
което е и причината трите ѝ поведения, тоест зоната на нечувствителност, несиметричните
периоди на изчакване и обработката на нулиран брояч, да могат да бъдат обсъждани по
измерени редове, а не по общи съображения.

Мащабирането в използваната среда е подлинейно и се изчерпва при 4 реплики, но
причината е известна и е в средата, не в услугата: 12 логически процесора, върху които
работят и генераторът, и посредникът, и репликите.

За твърдението „едно декларативно описание е достатъчно за двама доставчици“ проектът
дава частичен и отрицателен отговор, при това без да е разгръщал каквото и да било.
Общото описание на услугата е 82 реда, а двата платформено зависими модула са 297 и 178
реда. Тоест дели се \emph{описанието на услугата}, а не описанието на средата. Това е
резултат от преброяване на редовете в написаното описание, не от измерване при работещо
разгръщане, и е единственото, което може да се твърди преди прилагане.

\TODO{След действително разгръщане: доколко от разликата между двамата доставчици е
видима за потребителя на услугата, и колко от нея се вижда още в описанието на
инфраструктурата.}
